\section{Conclusions}

As transit bus operators face difficult decisions on designing their future fleets, decision support tools are of tremendous value. Designing a future fleet requires solving the \gls{vsfmp} which is known to be a NP-hard \gls{cop}. Solving the \gls{vsfmp} using an optimal solver requires evaluation of a set of integer trip-to-trip adjacency matrices (one for each bus type) and a set of supply-event=to-supply-event adjacency matrices (one for each supply port type). This complexity renders the problem effectively unsolvable for networks of realistic scale. Solving returns a single optimum which is useful for operational purposes but less so for strategic decision making. Strategic decision makers are better informed when they are able to see the trade-offs between objectives and how these manifest in fleet design.

This study alleviates this issue by presenting a flexible, efficient Pareto optimization algorithm for the \gls{vsfmp}. The \gls{nsga} framework utilized allows for a Pareto front to be revealed in N dimensions. The computational efficiency of the algorithm allows for reasonably quick iteration even using ordinary desktop or laptop computers which will be readily available to technical employees. The utility of the algorithm is demonstrated for medium and large bus networks in Section \ref{sec:case_study} efficiency and repeatability of the algorithm is demonstrated in Section \ref{sec:run_time}. Overall, the methodology presented in this paper, and implemented in the associated repository, provides transit agencies with a novel and usable tool to help them make their toughest decisions.

\subsection{Limitations and Future Work}

The methodology presented shares limitations with other metaheuristic methods and especially other \gls{nsga} based methods. Namely, it is not possible to know if a true optimum has been reached or if the solution arrived at is dictated by hyperparameters or even random seeding. Any \gls{nsga} method bears the specific issue that maintenance of field spread can only be guaranteed with impractically large population sizes. The method presented here mitigates this issue by allowing for a larger initial population and, thus, initial variety in solutions. Analysis of convergence across seeds shows that consistency is possible but hyperparameters may need to be tuned by network. Future work will include the development of objectives and constraints that match the real-life considerations of transit operators, more realistic models of both networks and equipment, and further work towards solution consistency including hyperparameter studies and alternate strategies such as population-migration.

\section{CRediT Statement}

\noindent\textbf{Aaron Rabinowitz:} Conceptualization, Methodology, Writing - Original Draft, Visualization, Software. \textbf{Jeremy Elvander:} Data Curation, Software, Writing - Original Draft. \textbf{Vaishnavi Karanam:} Conceptualization, Writing - Review \& Editing, Funding Acquisition, Project Administration. \textbf{Gil Tal:} Supervision, Funding Acquisition, Writing - Review and Editing.